\section{Diseño y Arquitectura de Software}

Se optó por una arquitectura cliente-servidor con microservicios. Este enfoque permite separar las funcionalidades por dominio, escalar cada servicio según su carga particular y desplegar cambios sin afectar al resto del sistema.

\subsection{Aplicación móvil}
\textbf{React Native} \autocite{reactnative}. Es el canal principal de uso para conductores y propietarios. Se eligió React Native para soportar iOS y Android con una sola base de código.

\subsection{Microservicios}
\textbf{Node.js} \autocite{nodejs} con \textbf{TypeScript} \autocite{typescript}. Implementan la lógica de negocio del sistema y exponen las APIs REST que consume la app móvil.

\subsection{Base de datos}
\textbf{PostgreSQL} \autocite{postgresql} con la extensión \textbf{PostGIS} \autocite{postgis}. Persiste los datos del sistema (usuarios, cocheras, reservas, pagos). PostGIS habilita las búsquedas de cocheras por proximidad geográfica.

\subsection{Caché}
\textbf{Redis} \autocite{redis}. Acelera el acceso a datos consultados frecuentemente: sesiones de usuario, tokens y disponibilidad de cocheras en tiempo real.

\subsection{Mensajería}
\textbf{RabbitMQ} \autocite{rabbitmq}. Permite que los microservicios se comuniquen de forma asincrónica mediante eventos. Por ejemplo, al procesar un pago, otros servicios reaccionan al evento sin bloquear el flujo principal.

\subsection{Observabilidad}
\textbf{Datadog} \autocite{datadog}. Centraliza métricas, logs y trazas distribuidas de cada microservicio para monitorear el sistema en producción.

\subsection{Mapas}
\textbf{Google Maps Platform} \autocite{googlemapsplatform}. Provee el mapa interactivo y las ubicaciones de las cocheras que ve el conductor.

\subsection{Pagos}
\textbf{MercadoPago} \autocite{mercadopago}. Procesa los pagos de las reservas. Se integra mediante SDK en el cliente para iniciar el cobro y webhooks en el backend para confirmar las transacciones.

\subsection{Notificaciones push}
\textbf{Firebase Cloud Messaging} \autocite{firebase}. Notifica en tiempo real a conductores y propietarios ante eventos relevantes, como la confirmación de una reserva o de un pago.

\subsection{Infraestructura}
\textbf{AWS} \autocite{aws} con cada servicio corriendo en un contenedor \textbf{Docker} \autocite{docker}. Sirve como base para desplegar y operar el sistema en producción.
